\scanpage{540}
\mspage{451}
% archive sheet number: 539

Dire que c'est l'identité, \deleted{\unsure{signifie que}} \inserted{\unsure{est|c'est}} exprimé par la condition $2^\circ)$ sur $(x_j,p_j)_{j\in J}$.

La condition $3^\circ$ est claire par construction, $p_j$ admettant une image isom. à $z$ par $\alpha_j$. D'autre part, on voit de suite que pour une flèche $\lambda:\inserted{j\mathrel{\relbar}j'}$ \unsure{ds|\deleted{de}} $J$, on a commutativité du

\diagram{\begin{tikzcd}[ampersand replacement=\&] x_j\arrow[r,dash,"{x_\lambda}"] \& x_{j'} \\ x_j^\natural\arrow[u,arrows={Hooks[left]-}]\arrow[r,dash,"{x_\lambda^\natural}"] \& x_{j'}^\natural\arrow[u,arrows={Hooks[left]->}] \\ z\arrow[u,"{\overline\alpha_j}"',"\sim"{rotate=90}]\arrow[r,dash,"{\underset{\sim}{\mathrm{id}_z}}"] \& z\arrow[u,"{\overline\alpha_{j'}}"',"\sim"{rotate=90}] \end{tikzcd}},

où $\overline\alpha_j,\overline\alpha_{j'}$ sont les iso. définis par $\alpha_j,\alpha_{j'}$ respectivement. Cela prouve que $x_\lambda^\natural$ est un iso.

Cela prouve que si $(x_i)_{i\in I}$ est ess. constant, il existe $J\mathrel{\relbar} I$ comme indiqué. L'inverse est également vrai, car OPS $J=I$, \unsure{puisque|\deleted{puis}que} $(x_i)_{i\in I}$ et $(x_j)_{j\in J}$ définissent des Ind-objets isomorphes. Alors le système inductif des $x_j^\natural$ a des morphismes de transition des iso, et comme $J$ est \inserted{filtrante donc} 1-connexe, on en déduit \unsure{facilemt|facilement} qu'il est isomorphe
